\section{Introduction}
\label{sec:intro}
% TODO(authors): confirm UAS vs ATM framing (Q-01). The motivation below is unchanged (UAS/AAM);
% the experiments use BlueSky transport aircraft and the text says so explicitly.
\lettrine{C}{ollision} avoidance among uncrewed aircraft systems (UAS) is a central problem
for the safe scaling of low-altitude airspace operations, including package delivery,
infrastructure inspection, and emerging UAS traffic management concepts. As operational
density increases, a fleet of cooperating vehicles must keep its members separated from one
another while each vehicle still completes its own mission.
% TODO(authors): deferred per decision 2026-11-25 (GAP-07/GAP-09/GAP-12) -- the original
% sentence also listed separation from static obstacles and non-cooperative intruders.
Advanced Air Mobility (AAM) operations impose a further constraint: the number of vehicles
sharing an airspace volume at any moment is not fixed at design time. Vehicles rotate out
for battery swaps, return to base for maintenance, are added as demand increases, or are
removed following mission completion. A coordination algorithm that must be retrained for
each fleet size $n$ is therefore difficult to operate.

Reinforcement learning offers a route to coordination policies that are not hand-tuned,
but most multi-agent reinforcement learning (MARL) methods share a structural limitation:
a model trained for a fixed fleet size $n$ does not transfer to a different fleet size
without architectural changes and retraining. In the centralized-training,
decentralized-execution (CTDE) paradigm, value-decomposition methods such as
VDN~\cite{vdn} and QMIX~\cite{qmix} train per-agent action-value functions through a
centralized mixing function. The QMIX mixing network has one input slot per agent and
hypernetworks whose output size is proportional to $n$, and the agent networks usually read
an observation vector formed by concatenating the features of all observed entities, whose
width also grows with $n$. Changing $n$ changes the shape of both networks.

Prior work has combined graph neural networks with CTDE, with different consequences for
this coupling. QMIX-GNN~\cite{qmix-gnn} enriches individual agent observations using a
graph attention network but passes the full set of $n$ augmented observations into the
centralized mixer, preserving the dependence on $n$. InforMARL~\cite{nayak2023scalable}
applies a mean-pooling operation over GNN-produced node embeddings to give the centralized
critic a fixed-size input, but this approach is restricted to actor-critic (PPO-family)
methods and cannot be applied to value decomposition architectures such as QMIX, since mean
pooling destroys the per-agent Q-value structure those methods require.
GraphMIX~\cite{naderializadeh2021graphconvolutionalvaluedecomposition} replaces QMIX's MLP
mixer with a GNN-based mixing module operating on a complete directed graph over all
agents; however, because the graph is complete and the input consists of one Q-value per
agent node, the mixing network remains structurally dependent on $n$ and requires
fine-tuning when $n$ changes.
% NOTE(authors): the three characterizations above are carried over unchanged from the
% extended abstract and should be checked against the cited papers.
TransfQMix~\cite{transfqmix} removes the coupling for value-based methods: it represents
observations and the state as matrices of entity features, processes them with transformers
in both the agent and the mixer, and generates the monotonic mixing weights from per-agent
embeddings, so that one set of parameters can be applied to teams of different sizes. Its
evaluation covers the Spread and StarCraft~II benchmarks.

This work adopts the TransfQMix principle and adapts it to cooperative air traffic conflict
resolution. The proposed method, GNN-QMIX, keeps the entity-matrix inputs, the decentralized
execution, and the generation of mixing weights from per-agent embeddings, and differs in
five respects. (i) The transformers are replaced by two graph neural networks: a graph
convolutional network (GCN)~\cite{gcn} in each agent and a graph attention network
(GAT)~\cite{gat} in the mixer. (ii) The agent does not learn a fully connected graph over
all entities; it operates on an explicit graph that contains only the vehicles inside its
sensor radius $r_{\text{detect}}$. (iii) There are no recurrent tokens in the agent or the
mixer; the bias and output-layer terms of the mixer are generated from fixed-size mean and
max readouts. (iv) The domain is horizontal conflict resolution between aircraft in the
BlueSky simulator~\cite{bluesky,bluesky-gym}. (v) One model is trained jointly on several
fleet sizes and evaluated zero-shot on larger ones.

The work is presented in two stages. Stage~1, GNN-QMIX, is learned end to end: the learned
policy issues the heading and speed command of every vehicle at every control step.
Stage~2, PID + GNN-QMIX, is a hybrid controller. A fixed proportional--integral--derivative
(PID) controller tracks each vehicle's waypoint through the same heading and speed commands,
and the learned conflict-resolution policy overrides it for a vehicle only while at least
one other vehicle is inside that vehicle's sensor radius. The same radius rule therefore
defines both the membership of the observation graph and the activation of the learned
policy.

\textbf{Theoretical claims.}
\begin{enumerate}
\item \emph{Fleet-size independence with monotonic mixing.} GNN-QMIX can be evaluated and
      trained at any fleet size with one set of parameters, and its mixer satisfies the
      QMIX monotonicity condition at every fleet size. This follows from three mechanisms:
      parameter sharing over graph nodes in the GCN agent and the GAT mixer, a first-layer
      mixing weight row that is generated separately for each agent from that agent's node
      embedding, and fixed-size mean and max readouts for the remaining mixing terms. The
      claim is stated and justified as Propositions~1 and~2 in Section~\ref{subsec:properties}.
\item \emph{Hybrid decomposition.} In the hybrid controller the waypoint-tracking layer is
      fixed, so the learner has to produce actions only for the vehicle-steps in which the
      radius trigger is active. We hypothesize that this reduces the effective
      state--action space that must be explored and improves sample efficiency relative to
      the end-to-end stage. This claim is tested empirically in the ablations of
      Section~\ref{subsec:results-ablations}; it is not proven here.
% TODO(authors): deferred per decision 2026-11-25 (GAP-07/GAP-09/GAP-12) -- theoretical claim 3
% (graceful degradation of coordination quality with communication latency \delta, mutual
% information bound) was removed with the latency model. Restore here when GAP-07 is resumed.
\end{enumerate}

\textbf{Contributions.}
\begin{enumerate}
\item A formulation of multi-aircraft conflict resolution in which each vehicle's
      observation and the global state are entity matrices, and the observation is a graph
      whose membership is set by a sensor radius (Section~\ref{subsec:entities}).
\item GNN-QMIX, a value-decomposition architecture with a GCN agent network and a monotonic
      GAT mixer, in which no parameter dimension depends on the number of vehicles
      (Sections~\ref{subsec:agent} and~\ref{subsec:mixer}).
\item A joint training scheme over several fleet sizes with one replay buffer per size and
      a summed loss, which needs no padding (Section~\ref{subsec:training}).
\item A hybrid PID + GNN-QMIX controller in which the sensor-radius rule switches control
      between a fixed waypoint-tracking law and the learned policy, together with an
      intervention count that measures how often the learned policy acts
      (Section~\ref{subsec:hybrid}).
\item An experimental protocol in BlueSky with joint training at $n\in\{3,6,9\}$, zero-shot
      evaluation at $n\in\{10,20,50\}$, QMIX, VDN, and IQL baselines, and separation metrics
      that are comparable across fleet sizes (Section~\ref{sec:setup}).
\end{enumerate}
% TODO(results): add one contribution item (or one closing paragraph) summarizing the measured
% outcome once Section~\ref{sec:results} is filled from summary_training.csv and
% summary_zero_shot.csv.
% TODO(authors): the extended abstract listed two further constructions that are no longer in
% scope and were removed from this list: (a) GNN + GRU modeling of the broadcast channel
% (deferred per decision 2026-11-25, GAP-07/GAP-09/GAP-12), and (b) supplying the PID action as
% an input to the learned policy (not decided for stage 2; see the TODO in
% Section~\ref{subsec:hybrid}). The emergence analysis was removed by author decision.
